\documentclass[12pt,a4paper]{article}
\usepackage[utf8]{inputenc}
\usepackage[english]{babel}
\usepackage{graphicx}
\usepackage{listings}
\usepackage{xcolor}
\usepackage{hyperref}
\usepackage{geometry}
\usepackage{microtype}
\usepackage{parskip}

\geometry{top=2.5cm, bottom=2.5cm, left=3cm, right=3cm}

\definecolor{javared}{rgb}{0.6,0,0}
\definecolor{javagreen}{rgb}{0.25,0.5,0.35}
\definecolor{javapurple}{rgb}{0.5,0,0.35}
\definecolor{javacomment}{rgb}{0.5,0.5,0.5}

\lstset{
    language=Java,
    basicstyle=\ttfamily\small,
    keywordstyle=\color{javapurple}\bfseries,
    stringstyle=\color{javared},
    commentstyle=\color{javacomment},
    numbers=left,
    numberstyle=\tiny\color{javacomment},
    frame=single,
    breaklines=true
}

\title{\textbf{Technical Report}\\
\large Dijkstra Routing Engine for the Paris Public Transport Network (IDFM)}
\author{Slims \\ M1 Software Engineering --- University Project}
\date{\today}

\begin{document}

\maketitle

\tableofcontents
\newpage

% ----------------------------------------------------------------
\section{Introduction}

This report describes the design and implementation of a data extraction and shortest-path routing system for the Paris Île-de-France Mobilités (IDFM) public transport network. The objective was to build a complete pipeline capable of downloading, processing, and modeling transport data, then applying Dijkstra's algorithm to compute optimal itineraries between any two stations.

The system is implemented in Java and exposed through both a command-line interface (CLI) and a JavaFX graphical interface.

% ----------------------------------------------------------------
\section{System Architecture}

The project follows a five-layer modular architecture designed for extensibility and separation of concerns:

\begin{itemize}
    \item \textbf{Presentation} --- CLI (\texttt{RouteDemo}) and GUI (\texttt{RouteFinderApp}, \texttt{RouteFinderController})
    \item \textbf{Application} --- Entry point and orchestration (\texttt{App})
    \item \textbf{Domain} --- Network modeling (\texttt{TransportGraph}, \texttt{Ligne}, \texttt{Station}, \texttt{Segment})
    \item \textbf{Infrastructure} --- Data ingestion and CSV I/O (\texttt{IDFMDataDownloader}, \texttt{CSVParser}, \texttt{CSVTools})
    \item \textbf{Algorithm} --- Graph and pathfinding (\texttt{Node}, \texttt{Dijkstra})
\end{itemize}

% ----------------------------------------------------------------
\section{Implementation}

\subsection{Maven Configuration}

The project uses Maven for dependency management. Key dependencies include:

\begin{lstlisting}[language=XML]
<dependency>
    <groupId>com.opencsv</groupId>
    <artifactId>opencsv</artifactId>
    <version>5.4</version>
</dependency>
<dependency>
    <groupId>org.json</groupId>
    <artifactId>json</artifactId>
    <version>20231013</version>
</dependency>
\end{lstlisting}

\subsection{Data Ingestion}

The extractor downloads two primary datasets from the IDFM open API:

\begin{itemize}
    \item \texttt{traces.csv} (~7.27 MB) --- line trajectories with GPS coordinates
    \item \texttt{stops.csv} (~12.15 MB) --- station names and positions
\end{itemize}

Travel durations between stops are estimated using the \textbf{Haversine formula}, which computes the great-circle distance between two GPS coordinates and applies an average speed assumption to derive travel time.

\subsection{Data Processing Pipeline}

The extraction process follows these sequential steps:

\begin{enumerate}
    \item Create the output directory
    \item Download CSV files from the IDFM API (with GZIP decompression support)
    \item Parse and process line trajectories (\texttt{TraceEntry} objects)
    \item Parse and match stop coordinates
    \item Clean invalid or unresolved entries via \texttt{cleanTraces()}
    \item Generate the structured \texttt{network.csv} via \texttt{CSVFormatter}
    \item Generate theoretical timetables via \texttt{ScheduleGenerator}
    \item Remove temporary files
\end{enumerate}

\subsection{Routing Algorithm}

The routing engine is built on \textbf{Dijkstra's algorithm} applied to a weighted undirected graph:

\begin{itemize}
    \item Each \textbf{node} represents a transit stop on a specific line (identified by \texttt{stopName\_lineId})
    \item Each \textbf{edge} encodes a travel duration in minutes
    \item \textbf{Transfer penalties} of 5 minutes are introduced between nodes of different lines sharing the same station name, modeling real-world interchange costs
\end{itemize}

This design allows the engine to naturally balance in-vehicle time against connection time --- a classic trade-off in transport network optimization.

\begin{lstlisting}
GrapheTransport graphe = new GrapheTransport();
graphe.loadFromCSVProvider(provider);

Station depart  = new Station("S1", "Station 1", 48.8566, 2.3522);
Station arrivee = new Station("S2", "Station 2", 48.8600, 2.3500);

Map<Station, Double> distances =
    Dijkstra.calculerChemin(graphe, depart, arrivee, true);
\end{lstlisting}

% ----------------------------------------------------------------
\section{Error Handling and Validation}

The system includes robust error handling at every stage:

\begin{itemize}
    \item Command-line argument validation
    \item Downloaded file integrity checks
    \item Detailed error messages with stack traces
    \item Graph validation: detection of isolated stations and invalid GPS coordinates via \texttt{GrapheTransport.validate()}
    \item Logging of key events and progress during data generation
\end{itemize}

% ----------------------------------------------------------------
\section{Generated Outputs}

\begin{itemize}
    \item \texttt{traces.csv} --- raw line trajectories (~7.27 MB)
    \item \texttt{stops.csv} --- raw stop data (~12.15 MB)
    \item \texttt{network.csv} --- structured network with durations and distances
    \item \texttt{schedules.csv} --- theoretical timetables (departures every 30 min, 7:00--9:00 AM)
\end{itemize}

% ----------------------------------------------------------------
\section{Build and Execution}

\textbf{Compile:}
\begin{lstlisting}[language=bash]
mvn clean compile
\end{lstlisting}

\textbf{Run the JavaFX interface:}
\begin{lstlisting}[language=bash]
mvn javafx:run
\end{lstlisting}

\textbf{Download and format IDFM data:}
\begin{lstlisting}[language=bash]
java -cp target/project.jar \
  fr.u_paris.gla.project.idfm.IDFMDataDownloader
\end{lstlisting}

% ----------------------------------------------------------------
\section{Conclusion}

This project demonstrates the design of a complete transport routing system, from raw API data ingestion to interactive shortest-path computation. The modular architecture ensures extensibility --- future improvements could include real-time delay integration, a REST API backend, and performance optimization for large-scale graph queries. The algorithmic challenges encountered here --- flow optimization in large networks, handling noisy real-world data, designing extensible systems --- are directly relevant to quantitative modeling in financial engineering contexts.

\end{document}